\documentclass[10pt]{article}

\usepackage[margin=0.82in]{geometry}
\usepackage[T1]{fontenc}
\usepackage[utf8]{inputenc}
\usepackage{lmodern}
\usepackage{microtype}
\usepackage{amsmath}
\usepackage{booktabs}
\usepackage{longtable}
\usepackage{tabularx}
\usepackage{array}
\usepackage{graphicx}
\usepackage{caption}
\usepackage{enumitem}
\usepackage{xurl}
\usepackage[numbers]{natbib}
\usepackage[hidelinks]{hyperref}

\graphicspath{{../results/figures/}}
\newcolumntype{Y}{>{\raggedright\arraybackslash}X}
\newcommand{\code}[1]{\texttt{\detokenize{#1}}}
\newcommand{\filepath}[1]{\path{#1}}
\newcommand{\fscore}{\ensuremath{F_1}}
\setlist{nosep,leftmargin=*}
\captionsetup{font=small}

\title{Text Classification Pipeline Forensics}
\author{AIAA3102 Final Project, Topic A}
\date{2026-07-21}

\begin{document}
\maketitle

\begin{abstract}
This project investigates a CPU-only classifier for the Disaster Tweets binary classification task, where target $1$ denotes a tweet describing a real disaster. The objective is not simply to maximize target-1 \fscore: it is to determine which pipeline changes have credible evidence, which errors they exchange, and which conclusions remain uncertain. All learned transformations are fitted on training data only; each decision is selected on dev, frozen, and then evaluated once on held-out data. The final dev-selected operating point is balanced Logistic Regression with $C=10$ and threshold $0.56$. It reaches held-out \fscore{} $0.754717$, compared with $0.731984$ for the required baseline. The reported gain is accompanied by explicit false-positive/false-negative trade-offs, surface-reliance stress tests, and a conservative data-quality audit that does not alter labels.
\end{abstract}

\section{Project Problem and Goal}

The project studies binary classification on the public Kaggle Disaster Tweets training data mirrored by UC Berkeley RISE \citep{kaggleDisasterTweets}. A tweet with target $1$ reports a disaster or related real-world event; target $0$ does not. The supplied split is fixed and is never regenerated: 4,567 training rows, 1,523 dev rows, and 1,523 held-out rows. The corresponding positive counts are 1,962, 655, and 654.

The central question is not whether a different score can be produced, but whether a score change has a defensible mechanism. The project therefore asks whether the required TF--IDF plus Logistic Regression baseline reproduces the course reference; whether normalization and feature-source changes reduce surface dependence rather than exploit shortcuts; which decision rule gives the strongest dev-selected precision-recall operating point; and which duplicate or error evidence warrants an audit action rather than a label change.

The course contract reports reference \fscore{} $0.757422$. The submitted baseline reaches held-out \fscore{} $0.731984$, a gap of $0.025438$. This gap is treated as a bounded forensic diagnosis rather than an invitation to choose whichever configuration is closest on held-out data.

\section{Methodology}

\subsection{Split, fitting, and freeze protocol}

Every vectorizer, categorical encoder, numeric scaler, and classifier is fitted only on the training partition. A dev command writes a pending decision; \code{freeze-ticket} copies that decision into \code{experiments/decisions.json}; only a frozen decision can run against held-out rows. \code{run-all} is the deterministic final regeneration command. Thus held-out data evaluates frozen configurations rather than choosing them.

\begin{table}[htbp]
\centering
\small
\begin{tabularx}{\textwidth}{l r Y Y}
\toprule
Stage & Rows available & Permitted action & Persisted evidence \\
\midrule
Train & 4,567 & Fit text and metadata transformations; fit model parameters & Fitted-ID checks and model configuration \\
Dev & 1,523 & Compare predefined candidates; select threshold, feature set, or model & Pending decision, dev predictions, candidate artifacts \\
Freeze & 0 additional & Promote one dev choice without re-fitting & \code{experiments/decisions.json} \\
Held-out & 1,523 & Score the frozen configuration once; inspect transitions & Held-out predictions, summary, confusion counts \\
\bottomrule
\end{tabularx}
\caption{Evaluation protocol.}
\label{tab:protocol}
\end{table}

All reported values are target-1 precision, recall, \fscore, and accuracy. For Tickets 2--5, \code{fixed_fp}, \code{fixed_fn}, \code{new_fp}, and \code{new_fn} are calculated against the same frozen Ticket 1 baseline, rather than against the immediately preceding ticket. This makes error movement comparable across the project.

The verified environment is Python 3.12.9 with NumPy 2.5.1 \citep{numpy}, pandas 3.0.3 \citep{pandas}, SciPy 1.18.0 \citep{scipy}, scikit-learn 1.9.0 \citep{scikitLearn}, and Matplotlib 3.11.0. The seed is 3102. The full environment is recorded in \code{results/environment.json}, while the final artifacts are reconstructed from prediction files by the validator.

\subsection{Ticket design}

\begin{table}[htbp]
\centering
\small
\begin{tabularx}{\textwidth}{c Y Y Y}
\toprule
Ticket & Hypothesis and controlled lever & Dev selection evidence & Frozen held-out role \\
\midrule
1 & A baseline mismatch is caused by unspecified implementation details; vary one factor per diagnostic probe. & Fixed 26-probe matrix; baseline itself remains assignment-driven. & Forensic dev/held-out delta agreement only; probes cannot replace the baseline. \\
2 & Stable URL replacement removes opaque URL-fragment reliance while retaining link presence. & Twelve one-lever normalization candidates. & Score the selected \code{url_replace} model and inspect transitions. \\
3 & Keyword, location, and shallow counts may be useful but brittle shortcut signals. & Seven feature-source configurations and metadata interventions. & Score the retained text-only model. \\
4 & A limited choice of regularization, class weighting, threshold, or LinearSVC can improve the operating point. & Fourteen Logistic Regression configurations plus one LinearSVC. & Score the frozen best dev configuration. \\
5 & Duplicates, conflicts, and confident model errors expose audit evidence but do not automatically prove a label error. & No label-correction experiment is selected. & Audit the frozen Ticket 4 model without changing labels or rows. \\
\bottomrule
\end{tabularx}
\caption{Controlled investigation design.}
\label{tab:ticket-design}
\end{table}

\section{Main Evidence and Results}

\subsection{Ticket 1: Baseline discrepancy diagnosis}

The frozen baseline uses lowercased word $(1,2)$-gram TF--IDF with sublinear term frequency, IDF and smooth IDF, a 40,000-feature cap, per-document L2 normalization, and \code{liblinear} Logistic Regression with L2 penalty, $C=1$, seed 3102, and threshold $0.5$. It obtains dev \fscore{} $0.734034$ and held-out \fscore{} $0.731984$. Its held-out precision is $0.777969$, recall is $0.691131$, and accuracy is $0.782666$.

The diagnostic matrix fixes every unnamed variable, including the split, training-only fitting, threshold, and baseline parameters. It changes exactly one named factor at a time. The full 26-probe artifact includes negative controls as well as improvements. It was fixed before diagnostic held-out scoring and is explicitly not a held-out model-selection table.

\begin{table}[htbp]
\centering
\scriptsize
\begin{tabularx}{\textwidth}{Y r r r r Y}
\toprule
One-factor probe & Dev \fscore{} & Held-out \fscore{} & Dev delta & Held-out delta & Interpretation boundary \\
\midrule
Submitted baseline & 0.734034 & 0.731984 & 0.000000 & 0.000000 & Required comparison anchor \\

$C=10$ & 0.750982 & 0.752137 & +0.016948 & +0.020153 & Larger $C$ weakens L2 regularization \\
$C=4$ & 0.745067 & 0.756078 & +0.011033 & +0.024095 & Diagnostic regularization response \\
Unigrams only & 0.738812 & 0.751220 & +0.004778 & +0.019236 & Removes phrase features; not selected \\
\code{min_df=3} & 0.741176 & 0.750419 & +0.007143 & +0.018435 & Removes many rare terms \\
\code{norm=None} & 0.738333 & 0.758621 & +0.004299 & +0.026637 & Preserves TF--IDF vector magnitude \\
\code{lbfgs} solver & 0.734034 & 0.731826 & +0.000000 & -0.000158 & Solver alone does not explain the gap \\
No IDF & 0.720949 & 0.723370 & -0.013085 & -0.008613 & Negative control \\
Preserve case & 0.717619 & 0.709106 & -0.016415 & -0.022878 & Negative control \\
L1 normalization & 0.316176 & 0.288582 & -0.417857 & -0.443402 & Strong negative control \\
\bottomrule
\end{tabularx}
\caption{Representative rows from the fixed 26-probe baseline diagnostic matrix.}
\label{tab:ticket1-probes}
\end{table}

The 26 dev/held-out \fscore{} deltas have Pearson $r=0.996813$ and Spearman $\rho=0.874915$. Figure~\ref{fig:ticket1-agreement} visualizes this descriptive agreement. Its left panel expands near-zero deltas so that modest positive changes can be inspected; the right panel retains the full range, including the severe L1-normalization decline. The probes share the same data, baseline, and model family, so this association is not treated as independent-sample inference, a p-value claim, or proof of causality.

\begin{figure}[htbp]
\centering
\includegraphics[width=0.96\textwidth]{ticket1_probe_delta_agreement.png}
\caption{Ticket 1 fixed-probe dev/held-out delta agreement. The left panel magnifies near-zero changes; the right panel preserves the full negative range.}
\label{fig:ticket1-agreement}
\end{figure}

Two representative probes were designated after dev comparison and before held-out scoring to make aggregate \fscore{} movement inspectable at the row level.

\begin{table}[htbp]
\centering
\small
\begin{tabularx}{\textwidth}{l r r r r Y}
\toprule
Probe & Fixed FP & Fixed FN & New FP & New FN & What the counts show \\
\midrule
$C=10$ & 12 & 38 & 32 & 6 & Recall-oriented movement: many recovered positives but a substantial FP cost. \\
\code{norm=None} & 53 & 36 & 26 & 26 & Broad movement in both directions; not a free improvement. \\
\bottomrule
\end{tabularx}
\caption{Post-freeze representative probe transitions.}
\label{tab:ticket1-transitions}
\end{table}

For example, $C=10$ recovers ID 244, a real tweet about a shooting or airplane accident, while newly flagging ID 117's conversational travel-delay wording. \code{norm=None} correctly rejects ID 303's non-event ``annihilated'' promotional wording but newly misses ID 302's terse positive ``Annihilated Abs'' text. The stable-ID rows are retained in \code{results/discrepancy_error_transitions.csv}; these examples illustrate prediction movement, not label corrections.

The baseline-reference gap remains unresolved because the reference vectorizer, classifier parameters, package versions, reference dev score, and reference predictions are unavailable. The evidence bounds plausible implementation causes but cannot identify a unique hidden reference configuration.

\subsection{Ticket 2: Text normalization lever}

Ticket 2 holds the Ticket 1 model, text features, split, and threshold fixed while changing one normalization decision at a time. The hypothesis is that a URL placeholder preserves the existence of a link without allowing opaque URL fragments to become fragile lexical features. No combinations of normalization changes are searched.

\begin{table}[htbp]
\centering
\small
\begin{tabular}{l r r r r c}
\toprule
Candidate & Dev precision & Dev recall & Dev \fscore{} & Delta from raw & Selected? \\
\midrule
\code{url_replace} & 0.810219 & 0.677863 & 0.738155 & +0.004121 & Yes \\
\code{hashtag_remove} & 0.775338 & 0.700763 & 0.736167 & +0.002133 & No \\
\code{url_remove} & 0.833333 & 0.656489 & 0.734415 & +0.000381 & No \\
Raw baseline & 0.780069 & 0.693130 & 0.734034 & +0.000000 & No \\
\code{mention_replace} & 0.780656 & 0.690076 & 0.732577 & -0.001457 & No \\
\code{mention_remove} & 0.774744 & 0.693130 & 0.731668 & -0.002366 & No \\
\code{preserve_case} & 0.758503 & 0.680916 & 0.717619 & -0.016415 & No \\
\bottomrule
\end{tabular}
\caption{Selected Ticket 2 dev candidates. Emoji replacement/removal, hashtag stripping, and punctuation removal/repetition exactly reproduce raw dev \fscore{}.}
\label{tab:ticket2-candidates}
\end{table}

Figure~\ref{fig:ticket2-evidence} shows all 12 candidates as a scaled dot plot rather than a zero-baseline bar chart, so the small but real dev differences are visible without exaggerating bar lengths. Its stress panel compares the raw baseline and the frozen URL-replacement model under the same synthetic text perturbations.

\begin{figure}[htbp]
\centering
\includegraphics[width=0.96\textwidth]{ticket2_normalization_and_stress.png}
\caption{Ticket 2 normalization candidates and surface-signal stress results on dev. Red denotes the frozen URL-replacement choice.}
\label{fig:ticket2-evidence}
\end{figure}

\begin{table}[htbp]
\centering
\scriptsize
\begin{tabularx}{\textwidth}{Y Y r r r r Y}
\toprule
Model evaluated on dev & Perturbation & \fscore{} before & \fscore{} after & Change & Flips & Reading \\
\midrule
Raw baseline & Replace URLs & 0.734034 & 0.714029 & -0.020005 & 125 & Raw model relies on opaque URL fragments. \\
URL-replacement model & Replace URLs & 0.738155 & 0.738155 & 0.000000 & 0 & Intended invariance is achieved. \\
URL-replacement model & Remove punctuation & 0.738155 & 0.720212 & -0.017943 & 70 & The model remains sensitive to a different surface transformation. \\
Raw baseline & Replace mentions & 0.734034 & 0.731826 & -0.002208 & 3 & Small residual mention sensitivity. \\
\bottomrule
\end{tabularx}
\caption{Ticket 2 dev-only perturbation evidence.}
\label{tab:ticket2-stress}
\end{table}

After the dev decision was frozen, \code{url_replace} reached held-out \fscore{} $0.739817$ and accuracy $0.794485$ (TN=765, FP=104, FN=209, TP=445). Relative to Ticket 1, it fixes 32 FPs and 8 FNs while introducing 7 FPs and 15 FNs. ID 353 is promotional ``World Annihilation'' content that becomes a fixed FP after URL fragments no longer dominate the score. ID 237 is a real airplane-accident post that becomes a new FN, demonstrating the recall cost. The conclusion is limited: URL replacement removes demonstrated URL-fragment dependence and slightly improves this dev-selected model; it does not make the classifier generally surface-invariant.

\subsection{Ticket 3: Feature and shortcut audit}

Ticket 3 tests whether supplied metadata or shallow counts add trustworthy information beyond tweet text. All candidates share the split, seed, classifier family, and threshold. Numeric features are tweet length, counts of words, URLs, mentions, hashtags, punctuation, and digits, uppercase ratio, and repeated punctuation. The selected text-only model is therefore an affirmative decision to reject shortcut additions, not an absence of experiments.

\begin{table}[htbp]
\centering
\small
\begin{tabular}{l r r r r c}
\toprule
Feature source & Dev precision & Dev recall & Dev \fscore{} & Delta from text & Selected? \\
\midrule
Text & 0.780069 & 0.693130 & 0.734034 & 0.000000 & Yes \\
Text + numeric & 0.786972 & 0.682443 & 0.730989 & -0.003045 & No \\
Text + keyword & 0.747126 & 0.694656 & 0.719937 & -0.014097 & No \\
Keyword + location & 0.675039 & 0.656489 & 0.665635 & -0.068399 & No \\
Keyword only & 0.679612 & 0.641221 & 0.659859 & -0.074175 & No \\
Numeric only & 0.616327 & 0.461069 & 0.527511 & -0.206523 & No \\
Location only & 0.579439 & 0.094656 & 0.162730 & -0.571304 & No \\
\bottomrule
\end{tabular}
\caption{Ticket 3 feature-source ablation on dev.}
\label{tab:ticket3-candidates}
\end{table}

The left panel of Figure~\ref{fig:ticket3-evidence} retains every dev candidate. The right panel only perturbs models that consume metadata, avoiding the misleading result that changing an unused field proves robustness.

\begin{figure}[htbp]
\centering
\includegraphics[width=0.96\textwidth]{ticket3_feature_and_stress.png}
\caption{Ticket 3 feature-source ablation and metadata-reliance stress results on dev.}
\label{fig:ticket3-evidence}
\end{figure}

\begin{table}[htbp]
\centering
\scriptsize
\begin{tabularx}{\textwidth}{Y Y r r r r Y}
\toprule
Consuming model & Intervention & \fscore{} before & \fscore{} after & Change & Flips & Evidence boundary \\
\midrule
Keyword only & Blank keyword & 0.659859 & 0.000000 & -0.659859 & 618 & Direct reliance on supplied keyword. \\
Keyword only & Shuffle keyword & 0.659859 & 0.434231 & -0.225627 & 699 & Depends on row-to-keyword pairing. \\
Text + keyword & Blank keyword & 0.719937 & 0.690009 & -0.029928 & 247 & Text model still changes when keyword is removed. \\
Text + keyword & Shuffle keyword & 0.719937 & 0.616242 & -0.103695 & 418 & Stronger disruption from mismatched keyword metadata. \\
Location only & Blank location & 0.162730 & 0.000000 & -0.162730 & 107 & Location is weak but entirely field-dependent. \\
\bottomrule
\end{tabularx}
\caption{Ticket 3 dev-only metadata interventions.}
\label{tab:ticket3-stress}
\end{table}

Text-only therefore remains the frozen Ticket 3 choice and produces the baseline held-out \fscore{} $0.731984$. This does not say that \code{keyword} is invalid task information; it is supplied context and contains real benchmark signal. It says that adding it to text lowers dev \fscore{} here and that models using it show strong synthetic dependence. Candidate transitions contain IDs 86 and 132 becoming fixed FPs and IDs 25 and 36 becoming new FPs; the exhaustive rows are in \code{results/error_transitions.csv}. Coefficients and perturbations are evidence of association and reliance, not causal semantics or deployment performance.

\subsection{Ticket 4: Decision rule and model search}

Ticket 4 keeps text features fixed and compares 14 Logistic Regression configurations---seven values of $C$ times \code{None} or \code{balanced} class weights---plus one LinearSVC. Each Logistic Regression configuration chooses its threshold from 0.20 to 0.80 on dev in 0.01 increments. LinearSVC uses its decision function at zero. This is a deliberately limited CPU-only search, not an unrestricted hyperparameter sweep.

Figure~\ref{fig:ticket4-grid} records the best dev \fscore{} for every Logistic Regression configuration. In this parameterization, larger $C$ means weaker L2 regularization. The selected configuration is balanced LR with $C=10$ and threshold 0.56. LinearSVC is retained as a genuine second-classifier comparison rather than omitted because it narrowly loses the dev competition.

\begin{figure}[htbp]
\centering
\includegraphics[width=0.76\textwidth]{ticket4_model_grid.png}
\caption{Ticket 4 dev \fscore{} for each Logistic Regression regularization and class-weight configuration.}
\label{fig:ticket4-grid}
\end{figure}

\begin{table}[htbp]
\centering
\small
\begin{tabularx}{\textwidth}{Y r r r r Y}
\toprule
Dev comparison & Threshold & Precision & Recall & \fscore{} & What is isolated \\
\midrule
Frozen Ticket 1 baseline & 0.50 & 0.780069 & 0.693130 & 0.734034 & Required baseline \\
Baseline threshold only & 0.46 & 0.750000 & 0.741985 & 0.745971 & Threshold effect alone \\
LR $C=1$, balanced & 0.51 & 0.741641 & 0.745038 & 0.743336 & Class weight near baseline regularization \\
LR $C=10$, unweighted & 0.52 & 0.787980 & 0.720611 & 0.752791 & Regularization with no class weighting \\
LinearSVC $C=1$ & 0.00 & 0.772947 & 0.732824 & 0.752351 & Second CPU classifier \\
LR $C=10$, balanced & 0.56 & 0.786070 & 0.723664 & 0.753577 & Selected combination \\
\bottomrule
\end{tabularx}
\caption{Dev ablation for the Ticket 4 decision.}
\label{tab:ticket4-ablation}
\end{table}

The ablation distinguishes the selected combination from any single explanation. Threshold tuning alone improves dev \fscore{} by 0.011937, but does not reach the selected combination. The selected $C=10$ balanced model trades some recall relative to lower thresholds for precision near the original baseline.

\begin{figure}[htbp]
\centering
\includegraphics[width=0.82\textwidth]{ticket4_dev_precision_recall.png}
\caption{Ticket 4 dev precision-recall curve with the frozen operating point marked at threshold 0.56.}
\label{fig:ticket4-pr}
\end{figure}

\begin{figure}[htbp]
\centering
\includegraphics[width=0.82\textwidth]{ticket4_dev_f1_threshold.png}
\caption{Ticket 4 dev \fscore{} across thresholds for the selected Logistic Regression configuration.}
\label{fig:ticket4-threshold}
\end{figure}

Figures~\ref{fig:ticket4-pr} and \ref{fig:ticket4-threshold} use dev scores only. They explain the selected operating point; they do not use held-out data to move the threshold after freezing. The frozen model reaches held-out \fscore{} $0.754717$, accuracy $0.795141$, precision $0.776699$, and recall $0.733945$ (TN=731, FP=138, FN=174, TP=480). Relative to Ticket 1, it fixes 17 FPs and 37 FNs, while introducing 26 FPs and 9 FNs. The held-out result supports moderate agreement with the dev decision, but remains one evaluation of a selection made on a single dev split.

\subsection{Ticket 5: Data-quality and error audit}

Ticket 5 does not retrain a label-corrected model. It searches for exact duplicate text, normalized duplicate keys, character 3--5-gram TF--IDF near duplicates, conflicting labels inside duplicate groups, and high-confidence held-out errors from the frozen Ticket 4 model. A finding is classified as \code{keep_but_flag}, \code{ambiguous}, or \code{reject_false_positive}; no held-out label or row is modified.

\begin{table}[htbp]
\centering
\scriptsize
\begin{tabularx}{\textwidth}{Y Y r r r Y}
\toprule
Evidence type & Disposition & Evidence rows & Unique IDs & Held-out rows & Interpretation \\
\midrule
Normalized duplicate & \code{keep_but_flag} & 1,019 & 1,019 & 202 & Repeated normalized text is a quality concern, not a label error by itself. \\
Conflicting-label duplicate & \code{ambiguous} & 266 & 266 & 42 & Duplicate conflict needs independent annotation. \\
Exact duplicate & \code{keep_but_flag} & 179 & 179 & 41 & Same-label duplicates are retained and flagged. \\
High-confidence FN & \code{ambiguous} & 75 & 75 & 75 & Model disagreement is not proof of a bad label. \\
High-confidence FP & \code{ambiguous} & 48 & 48 & 48 & Most apparent FPs remain review candidates. \\
Near duplicate & \code{keep_but_flag} & 73 & 73 & 0 & Consistent near duplicates indicate redundancy. \\
Near duplicate & \code{ambiguous} & 27 & 27 & 0 & Conflicting near duplicates need review. \\
High-confidence FP & \code{reject_false_positive} & 1 & 1 & 1 & Manual review supports the original label. \\
\bottomrule
\end{tabularx}
\caption{Ticket 5 audit evidence summary.}
\label{tab:ticket5-summary}
\end{table}

\begin{figure}[htbp]
\centering
\includegraphics[width=0.86\textwidth]{ticket5_audit_distribution.png}
\caption{Ticket 5 audit evidence counts by issue type and disposition.}
\label{fig:ticket5-audit}
\end{figure}

The audit has 1,688 evidence rows. IDs 59 (held-out) and 68 (train) are exact duplicate negative examples with cosine similarity 1.0, so both are \code{keep_but_flag}, not corrections. IDs 353 and 390 normalize to the same text but have opposite labels, so the audit assigns \code{ambiguous} rather than guessing which annotation is wrong. ID 365 is a high-confidence FN with score 0.1080 and belongs to a conflicting duplicate group; it remains a review case. ID 198 is a high-confidence FP with score 0.8524: its text discusses lifetime odds of dying in an airplane accident rather than a reported event. Semantic review supports its original negative label, so the candidate model finding is deliberately rejected. The audit identifies reproducible risks and review priorities, but supplies no independent evidence strong enough to justify a label correction.

\subsection{Cross-ticket held-out summary}

\begin{table}[htbp]
\centering
\small
\begin{tabular}{c l r r r c c}
\toprule
Ticket & Frozen model or decision & Dev \fscore{} & Held-out \fscore{} & Accuracy & Fixed FP/FN & New FP/FN \\
\midrule
1 & TF--IDF $(1,2)$ LR & 0.734034 & 0.731984 & 0.782666 & 0 / 0 & 0 / 0 \\
2 & URL replacement & 0.738155 & 0.739817 & 0.794485 & 32 / 8 & 7 / 15 \\
3 & Text only & 0.734034 & 0.731984 & 0.782666 & 0 / 0 & 0 / 0 \\
4 & Balanced LR, $C=10$, $t=0.56$ & 0.753577 & 0.754717 & 0.795141 & 17 / 37 & 26 / 9 \\
5 & Ticket 4 audit-only & 0.753577 & 0.754717 & 0.795141 & 17 / 37 & 26 / 9 \\
\bottomrule
\end{tabular}
\caption{Final frozen held-out results.}
\label{tab:cross-ticket}
\end{table}

\section{Case Analysis}

The aggregate metrics conceal different kinds of model behavior. Table~\ref{tab:cases} connects representative stable IDs to their ticket-specific evidence and limits.

\small
\begin{longtable}{p{0.06\textwidth} p{0.13\textwidth} p{0.23\textwidth} p{0.25\textwidth} p{0.16\textwidth}}
\caption{Representative stable-ID cases.}\label{tab:cases} \\
\toprule
Ticket & Stable IDs & Observation & Why the case matters & Boundary \\
\midrule
\endfirsthead
\toprule
Ticket & Stable IDs & Observation & Why the case matters & Boundary \\
\midrule
\endhead
1 & 244, 117 & $C=10$ recovers a real accident tweet but newly flags conversational accident wording. & \fscore{} gain contains both repaired FN and new FP movement. & A probe example is not a label correction or a new baseline. \\
1 & 303, 302 & \code{norm=None} rejects non-event ``annihilated'' wording but misses terse positive wording. & Different preprocessing changes different semantic failure modes. & Two examples do not prove a global semantic mechanism. \\
2 & 353, 237 & URL replacement fixes promotional ``World Annihilation'' content but misses a real airplane-accident post. & URL normalization reduces fragment reliance but can remove useful context. & The gain is small and selected on one dev split. \\
3 & 86, 132, 25, 36 & Shortcut candidates create both fixed and new FPs. & Metadata-based changes are not uniformly beneficial at row level. & Exhaustive transitions, not these examples alone, support the decision. \\
4 & 244, 117, 2528 & Final model recovers ID 244, creates an FP for ID 117, and misses traffic-collision ID 2528. & Higher \fscore{} still contains visible FP and FN costs. & \fscore{} is not a deployment-specific cost function. \\
5 & 59/68, 353/390, 198 & Same-label duplicates are flagged; conflicting duplicates stay ambiguous; a manual review rejects a false label-error finding. & Audit disposition separates redundancy, ambiguity, and correction evidence. & No independent annotator was available. \\
\bottomrule
\end{longtable}
\normalsize

The Ticket 4 operating point is preferable for the assignment's target-1 \fscore{} only because 37 recovered FNs outweigh 9 newly missed positives under that metric. It also creates 26 new FPs while fixing 17. This is why the report presents transitions, precision-recall curves, and examples alongside \fscore{} rather than treating the final score as a complete account.

\section{Difficulties and Solutions}

\begin{enumerate}
\item \textbf{The reference configuration was incomplete.} The contract provided a reference \fscore{} but not the vectorizer, classifier, package versions, reference dev score, or predictions. I froze the required baseline, tested a fixed one-factor probe matrix, and restricted held-out probe use to post-freeze forensic agreement rather than selecting the closest score.
\item \textbf{Preventing evaluation leakage required an executable protocol.} A narrative rule is insufficient if a CLI command can accidentally score held-out data. The pipeline uses pending dev decisions, an explicit freeze command, frozen decision records, train-fitted ID checks, and regression tests that reject held-out execution without a frozen decision.
\item \textbf{Aggregate \fscore{} did not explain error movement.} I added stable-ID prediction files, exhaustive transition categories, and artifact reconstruction. The validator recomputes \fscore{}, accuracy, and every transition count from predictions against the single Ticket 1 baseline.
\item \textbf{Initial shortcut perturbations could have been misleading.} Perturbing a metadata field on a text-only model would merely show that an unused field has no effect. The final stress design perturbs \code{keyword} and \code{location} only for models that consume them, making the observed changes interpretable.
\item \textbf{Data-quality evidence was easy to overclaim.} Duplicate text and high-confidence disagreement do not determine the correct label. The audit uses explicit dispositions and keeps every original label and held-out row, reserving \code{reject_false_positive} for the manually checked non-event ID 198.
\end{enumerate}

\section{AI Usage Declaration}

AI assistance was used for code review, comparing the implementation with the handout, identifying missing experimental controls, proposing focused artifact exports, drafting and reorganizing documentation, and checking wording. It was treated as a proposal mechanism rather than a source of evidence. Every numerical claim was verified against regenerated CSV artifacts; quoted stable-ID examples were checked against the source data or audit records; the dev/freeze/held-out boundary was inspected in code; and the report was rebuilt locally. Final verification ran \code{validate-data}, \code{run-all}, \code{validate-artifacts}, and 60 automated tests successfully. A concise, non-verbatim work record is retained in \code{logs/chat.md}.

\section{Discussion and Limitations}

The project has several important limitations. First, a single dev split selects normalization, features, regularization, class weighting, classifier family, and threshold. Small dev differences may be selection noise, particularly after multiple comparisons. Held-out scores validate frozen decisions but do not eliminate this uncertainty.

Second, \fscore{} treats precision and recall symmetrically. Threshold $0.56$ is an operating point optimized for the assignment metric, not a calibrated policy for a deployment setting with known false-positive and false-negative costs. The precision-recall and \fscore{}-threshold figures reveal the trade-off but cannot choose a real-world cost ratio.

Third, the Ticket 1 delta association is descriptive. The 26 probes are not independent experiments, and several settings alter related aspects of the same model. It supports reproducible directional behavior in this pipeline, but cannot prove a causal explanation of the undisclosed reference configuration.

Fourth, the URL and metadata perturbations are synthetic interventions. They establish sensitivity to transformations applied here, not robustness to all possible social-media shifts. Similarly, near-duplicate cosine thresholds are heuristic and high model confidence is not independent human annotation. Ticket 5 intentionally does not turn these cues into label changes.

\section{Conclusion}

The required baseline is deterministic and reproducible in the submitted environment, but it does not match the undocumented course reference configuration. The fixed Ticket 1 matrix narrows plausible sources of that gap without allowing held-out diagnostics to redefine the baseline. URL replacement gives a small dev-selected improvement and removes a measurable reliance on opaque URL fragments, while remaining sensitive to punctuation changes. Keyword and location metadata provide some benchmark signal but are weaker than text and visibly brittle under the interventions used here.

The strongest final dev-selected configuration is balanced Logistic Regression with $C=10$ and threshold $0.56$, reaching held-out \fscore{} $0.754717$. Its gain is real in the fixed split, but it has explicit error costs and remains subject to selection uncertainty. The data-quality audit adds a final constraint: duplicated or ambiguous evidence should be reported and reviewed, not converted automatically into changed labels. Overall, the project supports careful, reproducible model forensics rather than score-only optimization.

\section{References and Appendix}

\bibliographystyle{plainnat}
\bibliography{references}

\appendix
\section{Artifact Map}

\begin{longtable}{p{0.34\textwidth} p{0.60\textwidth}}
\toprule
Question & Primary machine-checkable evidence \\
\midrule
Split and environment & \filepath{data/split_indices.json}, \filepath{results/environment.json} \\
Frozen decisions & \filepath{experiments/decisions.json}, \filepath{experiments/pending_decisions.json} \\
Ticket 1 probes & \filepath{results/discrepancy_comparison.csv}, \filepath{results/discrepancy_association.json}, \filepath{results/discrepancy_error_transitions.csv} \\
Ticket 2 candidates and stress & \filepath{results/ticket2_normalization_comparison.csv}, \filepath{results/perturbation_stress.csv} \\
Ticket 3 feature and stress evidence & \filepath{results/ticket3_feature_comparison.csv}, \filepath{results/top_features.csv}, \filepath{results/perturbation_stress.csv} \\
Ticket 4 grid and operating point & \filepath{results/ticket4_model_grid.csv}, \filepath{results/decision_ablation.csv}, \filepath{results/ticket4_dev_decision_curve.csv}, \filepath{results/ticket4_dev_precision_recall_curve.csv} \\
Ticket 5 dispositions & \filepath{results/data_quality_audit.csv}, \filepath{results/ticket5_audit_summary.csv} \\
Final reconstruction & \filepath{predictions/heldout_predictions.csv}, \filepath{results/summary.csv}, \filepath{results/error_transitions.csv} \\
\bottomrule
\end{longtable}

\section{Reproduction}

From the repository root, run:

\begin{verbatim}
conda run -n aiaa3102 python -m pytest -q
conda run -n aiaa3102 python -m pipeline.cli validate-data
conda run -n aiaa3102 python -m pipeline.cli run-all
conda run -n aiaa3102 python -m pipeline.cli validate-artifacts
cd report && latexmk -pdf report.tex
\end{verbatim}

The first four commands regenerate and validate the machine-checkable artifacts. The final command compiles this report; the figures are referenced relative to the \code{report/} directory.

\end{document}